\begin{abstract}
As Enterprise AI systems increasingly rely on Retrieval-Augmented Generation (RAG) and Agentic architectures to perform multi-hop reasoning over organizational data, the dominant approach of independently retrieving and assembling context per query presents severe scalability limits. This paper investigates Continuous Context Computation (CCC)---a paradigm where enterprise information is continuously processed to maintain a reusable semantic representation, decoupling context assembly from query-time inference. We propose a comprehensive benchmark framework designed to evaluate CCC against conventional baselines (Dense RAG, Hybrid RAG with Reranking, and Long-Context LLMs). Our reproducible framework utilizes two public datasets (MS MARCO, HotpotQA) and three large-scale synthetic enterprise datasets (spanning 512,000 documents across 9 simulated platforms). We establish a rigorous 12-level question taxonomy and introduce novel metrics for Context Reuse Ratio and Context Compression. By isolating query-time computation from pre-computation, this research explicitly determines under what workloads, query patterns, and enterprise scales CCC provides measurable advantages, offering a formal economic model for break-even query volumes.
\end{abstract}
